% submission_info.tex (generated; edit src/submission_info.template.tex) -- what to enter on the submission server of the American Physical Society for the Regular Article
% "Higgs as a Composite, and Quantum Gravity" (Physical Review D).  PROJECT spacetimes, folder
% Revision_SM_Higgs/paper_prd.  Generated from src/submission_info.template.tex by scripts/build_paper.py,
% which inserts the abstract, the page counts and the table of the files from the build of the manuscript.
\documentclass[11pt]{article}
\usepackage[margin=2.5cm]{geometry}
\usepackage{amsmath,amssymb}
\usepackage{longtable}
\usepackage{array}
\usepackage{listings}
\usepackage[colorlinks=true,linkcolor=blue,urlcolor=blue]{hyperref}
\hypersetup{pdftitle={Submission information: Higgs as a Composite, and Quantum Gravity},pdfauthor={Patrick Lee Nash}}
\lstset{basicstyle=\ttfamily\footnotesize,breaklines=true,columns=fullflexible,frame=single,keepspaces=true,xleftmargin=4pt,xrightmargin=4pt}
\setlength{\parindent}{0pt}
\setlength{\parskip}{6pt}
\newcommand{\LL}{{\rm LL}}
\newcommand{\SF}{{\rm SF}}
\newcommand{\SH}{{\rm SH}}

\begin{document}

\begin{center}
{\Large\bf Submission information}\\[4pt]
{\large ``Higgs as a Composite, and Quantum Gravity''}\\[2pt]
Physical Review D, Regular Article\\[2pt]
PROJECT spacetimes, November 22, 2026
\end{center}

\section{What this document contains}

Everything that the submission server of the American Physical Society (APS) asks for, prepared from the manuscript: the
journal and the article type, the title, the byline, the abstract, the subject classification, the data availability
statement, the editorial information with the suggested referees, the cover letter, and the files to upload. The author
performs every step on the server himself; Sec.~\ref{sec:steps} lists the steps in order, and Sec.~\ref{sec:confirm} lists
the statements that only the author can confirm.

\section{Journal, article type and length}

\begin{itemize}
\item Journal: Physical Review D. Article type: Regular Article (not a Letter, Comment or Reply; no earlier paper needs to be
named).
\item Length: the manuscript has 14 pages in the two-column format of REVTeX 4.2, the Supplemental Material
30 pages, the cover letter 2 pages. Physical Review D sets no limit to the length of Regular Articles
(the appropriate length depends on the information presented); the word limits of Letters and Comments do not apply.
\item Abstract: 417 words (a formula counted as one word); the APS guidance asks for less than 500 words.
\end{itemize}

\section{Title and byline}
\label{sec:title}

\begin{itemize}
\item Title: Higgs as a Composite, and Quantum Gravity
\item Author: Patrick Lee Nash, Ph.D. (the manuscript prints the name as the author asked; published Physical Review
bylines show names without academic degrees, so the production editors may print ``Patrick Lee Nash'').
\item Affiliation: PROJECT spacetimes. APS asks for the institution where the research was done, usually with the city,
the state and the country; the author completes it on the server (and in \texttt{\textbackslash affiliation} of the
source, if he wishes) with the address he wants printed.
\item Co-authors: none. ORCID: the server asks the author to link his ORCID iD.
\end{itemize}

\section{Abstract}
\label{sec:abstract}

As it appears in the manuscript:

\begin{quote}
Can the Higgs field be a condensate of pairs of fermions, as the state of a superconductor is a condensate of pairs of electrons? We study this question in a model built on an algebraic structure, the Lorentz super-loop: the bi-octonions of unit norm, a Moufang loop (a set with a product that has an identity and inverses but is not associative). It contains the Lorentz group $SL(2,\mathbb C)$ and eight spinor charges, which are carried by an isodoublet $\Psi$ of Dirac spinors. We first show, with worked examples, how these objects arise, and why pairs of two members of $\Psi$ form isospin triplets, so that a condensate of $\Psi$ alone cannot be the Higgs doublet of the Standard Model. We then give $\Psi$ the hypercharge $\frac14$, add two singlet Dirac fields $\chi$ and $\chi'$ with the hypercharges $\frac14$ and $-\frac34$, and let a contact interaction bind a member of $\Psi$ to a singlet in a pair of spin zero. The two pair fields form a bidoublet of $SU(2)_L\times SU(2)_R$ (a $2\times2$ matrix of fields on which the two groups act from the left and from the right) with the quantum numbers of two Higgs doublets. In five mean-field methods (Hubbard-Stratonovich with three decouplings, Kohn-Sham, and Cornwall-Jackiw-Tomboulis) the ground state is the custodial vacuum, which keeps the symmetry that protects the ratio of the $W$ and $Z$ masses: it breaks $SU(2)_L\times U(1)_Y$ to the $U(1)$ of the electric charge, keeps the photon massless and gives $\rho=1$. Setting its scale to $v=246.22$~GeV puts the momentum cutoff at 2.2\mbox{--}2.4~TeV. The amplitude mode, the oscillation of the size of the condensate and the candidate for the Higgs boson, is a composite scalar bound below the threshold of two quasi-particles. Where only the pairing acts, the linear response has the closed form of the theory of Bardeen, Cooper and Schrieffer and places the mode exactly at the pair self-energy (the gap parameter of that theory): $M_H/v=7.967$, 16 times the measured $m_H/v=0.5082$, and over a scan of the coupling from 1.02 to 3 times its critical value, with two masses of the fields, the ratio is at least \ensuremath{4.892}. A custodial triplet of massive scalars and a Goldstone boson of the particle number come with it. The superfluid therefore has the quantum numbers of the Higgs field and gives the $W$ and $Z$ bosons their masses, but in the mean field its amplitude mode is far too heavy. We outline how gravitation will enter through E.~A.~Lord's gauge theory of the local change of spinor representation.
\end{quote}

The LaTeX source of the same text, to paste into the abstract field of the form:

\begin{lstlisting}
Can the Higgs field be a condensate of pairs of fermions, as the state of a superconductor is a condensate of pairs of electrons? We study this question in a model built on an algebraic structure, the Lorentz super-loop: the bi-octonions of unit norm, a Moufang loop (a set with a product that has an identity and inverses but is not associative). It contains the Lorentz group $SL(2,\mathbb C)$ and eight spinor charges, which are carried by an isodoublet $\Psi$ of Dirac spinors. We first show, with worked examples, how these objects arise, and why pairs of two members of $\Psi$ form isospin triplets, so that a condensate of $\Psi$ alone cannot be the Higgs doublet of the Standard Model. We then give $\Psi$ the hypercharge $\frac14$, add two singlet Dirac fields $\chi$ and $\chi'$ with the hypercharges $\frac14$ and $-\frac34$, and let a contact interaction bind a member of $\Psi$ to a singlet in a pair of spin zero. The two pair fields form a bidoublet of $SU(2)_L\times SU(2)_R$ (a $2\times2$ matrix of fields on which the two groups act from the left and from the right) with the quantum numbers of two Higgs doublets. In five mean-field methods (Hubbard-Stratonovich with three decouplings, Kohn-Sham, and Cornwall-Jackiw-Tomboulis) the ground state is the custodial vacuum, which keeps the symmetry that protects the ratio of the $W$ and $Z$ masses: it breaks $SU(2)_L\times U(1)_Y$ to the $U(1)$ of the electric charge, keeps the photon massless and gives $\rho=1$. Setting its scale to $v=246.22$~GeV puts the momentum cutoff at 2.2\mbox{--}2.4~TeV. The amplitude mode, the oscillation of the size of the condensate and the candidate for the Higgs boson, is a composite scalar bound below the threshold of two quasi-particles. Where only the pairing acts, the linear response has the closed form of the theory of Bardeen, Cooper and Schrieffer and places the mode exactly at the pair self-energy (the gap parameter of that theory): $M_H/v=7.967$, 16 times the measured $m_H/v=0.5082$, and over a scan of the coupling from 1.02 to 3 times its critical value, with two masses of the fields, the ratio is at least \ensuremath{4.892}. A custodial triplet of massive scalars and a Goldstone boson of the particle number come with it. The superfluid therefore has the quantum numbers of the Higgs field and gives the $W$ and $Z$ bosons their masses, but in the mean field its amplitude mode is far too heavy. We outline how gravitation will enter through E.~A.~Lord's gauge theory of the local change of spinor representation.
\end{lstlisting}

\section{Subject classification (PhySH)}
\label{sec:physh}

The Physics Subject Headings are chosen by searching in the classification field of the form. The following concepts
were found in PhySH under \emph{Research Areas} $\to$ \emph{Hypothetical particle physics models} (discipline Particles
\& Fields): \emph{Composite models}, \emph{Dynamical symmetry breaking models}, \emph{Extensions of Higgs sector}. Suggested:
\begin{itemize}
\item primary concept (to be starred): \emph{Composite models};
\item further concepts: \emph{Dynamical symmetry breaking models}; \emph{Extensions of Higgs sector} (the superfluid has
the quantum numbers of two Higgs doublets);
\item further concepts to search for (their exact names in PhySH were not checked): custodial symmetry; electroweak
symmetry breaking; spontaneous symmetry breaking; collective modes; mean field theory; superfluidity; quantum gravity.
\end{itemize}
These names are those of the classification as it was found for the earlier version of this manuscript; the author
confirms them in the form.

\section{Data availability statement}
\label{sec:das}

Physical Review D publishes a data availability statement with every article, and the server asks for it with its
citations. The statement of the manuscript, with the citation written out:

\begin{quote}
The data that support the findings of this article are openly available in the repository of PROJECT spacetimes,
\url{https://github.com/once-ere/spacetimes}, folder \texttt{Revision\_SM\_Higgs}: the data files (\texttt{data/}), the
Rust, Python and Wolfram Language programs that produce and check them (\texttt{rust/}, \texttt{python/},
\texttt{wolfram/}), the reports of their checks, the notebook, and the scripts that build this article from the data
(\texttt{paper\_prd/}). The Supplemental Material lists the files with their SHA-256 checksums.
\end{quote}

A repository can change after publication. A frozen copy with a persistent identifier makes the statement permanent: for
example a release of the repository archived by a service that assigns a DOI (Zenodo archives releases of GitHub
repositories), whose DOI is then added to the statement and to the reference of the repository.

\section{Editorial information}
\label{sec:editorial}

\paragraph{Comments for the editor} (a proposal):
\begin{quote}
The Supplemental Material (file \texttt{higgs\_composite\_qg\_SM.pdf}) contains the derivations, the complete tables and
every numerical check with its deviation and tolerance; the programs and data are in a public repository (see the data
availability statement). This is not a joint submission. The manuscript builds on an unpublished manuscript of the author,
dated November 16, 2026, which can be provided on request; everything that the manuscript uses of it is restated and
derived in its Sec.~II, so that the manuscript can be read on its own.
\end{quote}

\paragraph{Suggested referees} (by field; the manuscript cites work of each). The editors choose the referees; the form
asks for names, institutions and e-mail addresses, which the author takes from the institutional web pages of the persons
he selects, after checking that none of them has a conflict of interest with him (a collaboration, a common institution,
a personal relation):
\begin{itemize}
\item Christopher T. Hill: dynamical symmetry breaking, top-quark condensation and the compositeness condition (W.~A.~Bardeen,
C.~T.~Hill and M.~Lindner, Phys.\ Rev.\ D 41, 1647 (1990); C.~T.~Hill and E.~H.~Simmons, Phys.\ Rep.\ 381, 235 (2003));
\item Roberto Contino: the composite Higgs boson as a pseudo-Goldstone boson (K.~Agashe, R.~Contino and A.~Pomarol, Nucl.\
Phys.\ B 719, 165 (2005));
\item Chandra M. Varma: the amplitude (Higgs) modes of condensates of fermion pairs (P.~B.~Littlewood and C.~M.~Varma,
Phys.\ Rev.\ Lett.\ 47, 811 (1981); D.~Pekker and C.~M.~Varma, Annu.\ Rev.\ Condens.\ Matter Phys.\ 6, 269 (2015));
\item Jo\~ao P. Silva: models with two Higgs doublets and their custodial symmetry (G.~C.~Branco \emph{et al.}, Phys.\
Rep.\ 516, 1 (2012)).
\end{itemize}

\paragraph{Joint submissions, excluded referees:} none.

\section{Cover letter}

The file \texttt{cover\_letter.pdf} (2 pages) gives the context of the manuscript (the earlier manuscript
on the Lorentz super-loop, why its superfluid could not be the Higgs field, and the self-contained Sec.~II that restates
it), the three main findings with their numbers, the reasons for Physical Review D, a statement on
the title and the gravitational sector (a program, not computed), the reproducibility of the numbers, the submission
history (a first submission, not under consideration elsewhere) and a pointer to the suggested referees.

\section{Files to upload}
\label{sec:files}

The folder \texttt{paper\_prd/submission/} holds exactly these files. Its LaTeX source refers to the figures without a
folder, because the server keeps the uploaded files in one folder; the build compiles a copy of the folder alone and checks
that it gives the same number of pages without an undefined reference. The server builds its own PDF from the source,
which the author compares with \texttt{higgs\_composite\_qg.pdf} before submitting.

\begingroup\small\setlength{\tabcolsep}{4pt}
\begin{longtable}{>{\raggedright\arraybackslash}p{0.31\textwidth}>{\raggedright\arraybackslash}p{0.31\textwidth}r>{\raggedright\arraybackslash}p{0.2\textwidth}}
\hline\hline file & content & kB & SHA-256 (first 16 digits) \\ \hline
\endfirsthead
\hline file & content & kB & SHA-256 (first 16 digits) \\ \hline
\endhead
\hline\hline
\endlastfoot
\texttt{higgs\_composite\_qg.pdf} & the manuscript (PDF) & 480 & \texttt{9bc5050bb86f3090} \\
\texttt{higgs\_composite\_qg.tex} & the manuscript (LaTeX source, REVTeX 4.2) & 80 & \texttt{06a4df0d487d13c8} \\
\texttt{higgs\_composite\_qg.bbl} & the bibliography of the source (from BibTeX with apsrev4-2.bst) & 26 & \texttt{01b06154a457703c} \\
\texttt{fig\_prd\_sh\_potential.pdf} & Figure 1 & 19 & \texttt{99c334598e4a3194} \\
\texttt{fig\_prd\_sh\_scan.pdf} & Figure 2 & 17 & \texttt{9d2ae12b30f5d189} \\
\texttt{fig\_prd\_sh\_modes.pdf} & Figure 3 & 19 & \texttt{df4627138882083d} \\
\texttt{higgs\_composite\_qg\_SM.pdf} & the Supplemental Material (PDF) & 580 & \texttt{4ab6551bd04407fa} \\
\texttt{cover\_letter.pdf} & the cover letter (PDF) & 87 & \texttt{b04bd6724fd309c6} \\
\end{longtable}
\endgroup

\section{The steps on the submission server}
\label{sec:steps}

\begin{enumerate}
\item Log in to the APS submission server and choose ``Create a New Submission''.
\item Check or update the contact information at the top of the first page.
\item Select the journal Physical Review D and confirm the journal-specific guidelines.
\item Select the article type Regular Article.
\item Enter the title and the abstract (Secs.~\ref{sec:title} and \ref{sec:abstract}).
\item Co-authors: none; link the ORCID iD.
\item Subject classification: search the PhySH concepts of Sec.~\ref{sec:physh} and star the primary concept.
\item Open Science information: enter the data availability statement of Sec.~\ref{sec:das}.
\item Editorial information: the comments for the editor and the suggested referees of Sec.~\ref{sec:editorial}.
\item Upload the cover letter and the files of Sec.~\ref{sec:files} (the Supplemental Material as Supplemental Material).
\item Compare the PDF built by the server with \texttt{higgs\_composite\_qg.pdf}, then submit.
\end{enumerate}

\section{Statements that only the author can confirm}
\label{sec:confirm}

\begin{itemize}
\item The submission history of the cover letter: the manuscript has not been published or submitted elsewhere.
\item The status of the earlier manuscript on the Lorentz super-loop (unpublished, dated November 16, 2026), which the
cover letter and the comments for the editor state.
\item The absence of conflicts of interest with the suggested referees.
\item The affiliation and the address to be printed, and the form of the name.
\item The funding information that the server asks for (the manuscript names no funding).
\item The title: the manuscript and the cover letter say that the gravitational sector is a program and is not computed;
the editors or the referees may ask whether the title should say so.
\end{itemize}

\end{document}
